% =============================================================================
%  PAPER 2 — T-FIRST PROGRAMME
%  "The Delta-Gain Closes the Leray-Prodi-Serrin Gap:
%   A Conditional Regularity Result for Navier-Stokes"
%
%  Target journal: Communications in Mathematical Physics
%  Authors: Pratik Menon, Claude (Anthropic)
%
%  STATUS: DRAFT — Not for citation or distribution
% =============================================================================

\documentclass[12pt]{article}

%----------------------------------------------------------------------
%  Packages
%----------------------------------------------------------------------
\usepackage{amsmath,amssymb,amsthm,amsfonts}
\usepackage{mathtools}
\usepackage{geometry}
\usepackage{hyperref}
\usepackage{enumitem}
\usepackage{booktabs}
\usepackage{array}
\usepackage{microtype}
\usepackage{xcolor}
\usepackage{cleveref}

\geometry{
  top=1.15in, bottom=1.15in,
  left=1.25in, right=1.25in,
  marginpar=0in
}

\hypersetup{
  colorlinks=true,
  linkcolor=blue!60!black,
  citecolor=blue!60!black,
  urlcolor=blue!60!black
}

%----------------------------------------------------------------------
%  Theorem environments
%----------------------------------------------------------------------
\theoremstyle{plain}
\newtheorem{theorem}{Theorem}[section]
\newtheorem{lemma}[theorem]{Lemma}
\newtheorem{proposition}[theorem]{Proposition}
\newtheorem{corollary}[theorem]{Corollary}

\theoremstyle{definition}
\newtheorem{definition}[theorem]{Definition}
\newtheorem{assumption}[theorem]{Assumption}
\newtheorem{example}[theorem]{Example}

\theoremstyle{remark}
\newtheorem{remark}[theorem]{Remark}

%----------------------------------------------------------------------
%  Custom commands
%----------------------------------------------------------------------
\newcommand{\R}{\mathbb{R}}
\newcommand{\T}{\mathbb{T}}
\newcommand{\N}{\mathbb{N}}
\newcommand{\Z}{\mathbb{Z}}
\newcommand{\C}{\mathbb{C}}
\newcommand{\Lp}[1]{L^{#1}}
\newcommand{\Hs}[1]{H^{#1}}
\newcommand{\Ws}[2]{W^{#1,#2}}
\newcommand{\norm}[1]{\left\|#1\right\|}
\newcommand{\abs}[1]{\left|#1\right|}
\newcommand{\inner}[2]{\left\langle #1,\, #2\right\rangle}
\newcommand{\pa}{\partial}
\newcommand{\eps}{\varepsilon}
\newcommand{\grad}{\nabla}
\newcommand{\Div}{\operatorname{div}}
\newcommand{\curl}{\operatorname{curl}}
\newcommand{\lap}{\Delta}
\newcommand{\nt}[1]{\nu\abs{\grad #1}^2}    % Calderón-Zygmund source
\newcommand{\LPS}{\text{LPS}}
\newcommand{\NS}{\text{NS}}
\newcommand{\supp}{\operatorname{supp}}
\newcommand{\diam}{\operatorname{diam}}
\newcommand{\sgn}{\operatorname{sgn}}
\newcommand{\openinterval}[2]{(#1,\,#2)}
\newcommand{\closeinterval}[2]{[#1,\,#2]}
\newcommand{\half}{\tfrac{1}{2}}
\newcommand{\tT}{[0,T]}

% Draft watermark colour
\newcommand{\draftcolor}{\color{red!70!black}}

%----------------------------------------------------------------------
%  Title block
%----------------------------------------------------------------------
\title{%
  {\draftcolor\textbf{[DRAFT --- Not for distribution]}}\\[6pt]
  \textbf{The Delta-Gain Closes the Leray-Prodi-Serrin Gap:\\
          A Conditional Regularity Result for the\\
          3D Navier-Stokes Equations}
}

\author{
  Pratik Menon\\[2pt]
  \small Independent Researcher\\[2pt]
  \small \texttt{pmenonn@gmail.com}
  \and
  Claude (Anthropic)\\[2pt]
  \small Anthropic\\[2pt]
  \small \texttt{claude@anthropic.com}
}

\date{April 2026}

%----------------------------------------------------------------------
\begin{document}
%----------------------------------------------------------------------

\maketitle
\thispagestyle{empty}

\begin{center}
  {\draftcolor\textbf{DRAFT — April 2026 — T-First Programme, Paper 2 of 3}}
\end{center}

\vspace{4pt}

\begin{abstract}
We prove a conditional regularity theorem for the three-dimensional incompressible
Navier-Stokes equations: if an auxiliary \emph{scalar} field $\theta$ satisfying
\[
  \pa_t\theta + u\cdot\grad\theta = \nu\,\lap\theta + \nu\abs{\grad u}^2,
  \qquad \theta(\cdot,0)=0,
\]
belongs to $L^2\bigl(0,T;\,H^{1+\delta}(\R^3)\bigr)$ for some $\delta>0$, then every
Leray-Hopf weak solution $u$ is smooth on $(0,T]$.
The argument proceeds through an explicit Sobolev embedding chain that places $u$
strictly inside the Prodi-Serrin class $L^p_t L^q_x$ with
$2/p+3/q<1$, yielding regularity by the classical Serrin criterion.
The threshold is sharp at the level of scaling: $\delta=0$ corresponds exactly
to the critical exponent pair $(p,q)=(\infty,3)$, which is the open endpoint of
the Leray-Prodi-Serrin gap.
We verify numerically, using spectral simulations in both two and three dimensions
with Wang et al.\ adversarial initial conditions, that $\delta\geq 0.50$
is observed for every initial condition tested, including the critical counter-flow
profile $\lambda=0.4703$ and the adversarial minimum $\lambda=0.05$.
The problem of proving $\delta>0$ unconditionally (Claim~A) is addressed in
Paper~3 of this programme.
\end{abstract}

\vspace{8pt}
\noindent\textbf{Keywords.}
Navier-Stokes equations, Leray-Hopf weak solutions, Prodi-Serrin regularity,
auxiliary scalar, Calderón-Zygmund estimates, Sobolev embedding, regularity gap.

\vspace{8pt}
\noindent\textbf{MSC 2020.}
35Q30, 35B65, 76D05.

\tableofcontents

\newpage

%======================================================================
\section{Introduction}
%======================================================================

\subsection{Background and motivation}

The Clay Millennium Problem on Navier-Stokes \cite{Fefferman2000} asks whether every
smooth, finite-energy initial datum in $\R^3$ gives rise to a global smooth solution
of the incompressible Navier-Stokes equations
\begin{equation}\label{eq:NS}
  \pa_t u + (u\cdot\grad)u = -\grad p + \nu\,\lap u,
  \qquad \Div u = 0,
\end{equation}
with $u(\cdot,0)=u_0\in C^\infty_c(\R^3)$, $\Div u_0=0$, and $\nu>0$ fixed.
Despite decades of effort, the answer remains unknown.
The deepest unconditional result is Leray's theorem \cite{Leray1934}: global
\emph{weak} solutions (Leray-Hopf solutions) exist, but may fail to be unique or smooth.

The principal obstruction to proving global smoothness is the \emph{supercritical}
scaling of the equations: under the rescaling $(u,p,t,x)\mapsto
(\lambda u,\lambda^2 p,\lambda^{-2}t,\lambda^{-1}x)$ the energy $\norm{u}_{L^2}^2$
scales as $\lambda$, not as a conserved quantity.
This gap between the energy level and the critical scaling is the fundamental
difficulty.

\subsubsection*{The Leray-Prodi-Serrin theorem}

The most powerful conditional regularity result, due independently to Prodi
\cite{Prodi1959} and Serrin \cite{Serrin1962} and extended by Ladyzhenskaya
\cite{Ladyzhenskaya1967}, states the following.

\begin{theorem}[Prodi-Serrin-Ladyzhenskaya, 1959/1962/1967]\label{thm:LPS}
  Let $u$ be a Leray-Hopf weak solution of \eqref{eq:NS} on $\R^3\times(0,T)$.
  If
  \begin{equation}\label{eq:serrin-condition}
    u\in L^p\bigl(0,T;\,L^q(\R^3)\bigr)
    \quad\text{with}\quad
    \frac{2}{p}+\frac{3}{q}\leq 1,\quad q\geq 3,
  \end{equation}
  then $u$ is smooth on $(0,T]$.
\end{theorem}

The condition \eqref{eq:serrin-condition} defines the \emph{Prodi-Serrin class}.
The \emph{Leray-Prodi-Serrin (LPS) gap} refers to the range of pairs $(p,q)$ for
which neither \eqref{eq:serrin-condition} holds unconditionally nor a counterexample
is known.
The critical endpoint $(p,q)=(\infty,3)$ — corresponding to the borderline condition
$u\in L^\infty_t L^3_x$ — is especially delicate.
Escauriaza, Seregin, and \v{S}ver\'ak \cite{ESS2003} proved that a solution
\emph{cannot} blow up if $\norm{u(t,\cdot)}_{L^3}\to\infty$ as $t\nearrow T^*$;
equivalently, if $u\in L^\infty(0,T;L^3)$ then the solution extends smoothly.
However, the full open endpoint $u\in L^\infty L^3$ remains, in its most general
formulation, unresolved at the level of the Prize problem.

\subsubsection*{The T-First programme}

The T-First programme attacks the Prize problem through an \emph{auxiliary scalar}
field $\theta$ that encodes the enstrophy production via
$\nu\abs{\grad u}^2$.
The central thesis of the programme, articulated in its primary document
\cite{TFirstPRD2026}, is:

\begin{center}
  \emph{Any $H^{1+\delta}$ regularity gain in $\theta$ — however small — suffices
  to close the LPS gap.}
\end{center}

This paper makes that statement precise and proves it rigorously.
The unconditional claim that $\delta>0$ always holds (``Claim~A'') is the subject
of Paper~3 \cite{TFirstPaper3}.

\subsection{The auxiliary scalar}

Define the auxiliary scalar $\theta:\R^3\times[0,T]\to\R$ by
\begin{equation}\label{eq:theta}
  \pa_t\theta + u\cdot\grad\theta = \nu\,\lap\theta + \nu\abs{\grad u}^2,
  \qquad \theta(\cdot,0) = 0.
\end{equation}
The source term $S_\theta := \nu\abs{\grad u}^2$ is precisely the enstrophy
production density (up to a constant).
Equation \eqref{eq:theta} is an advection-diffusion equation for $\theta$ driven
by this source; the scalar $\theta$ is \emph{slaved} to $u$ in the sense that once
$u$ is given, $\theta$ is determined.

The source structure $S_\theta\in L^1_{t,x}$ (which follows from the Leray energy
inequality) already places $\theta$ in $L^1_{t,x}$, but the critical question is
whether $\theta$ gains additional Sobolev regularity.
The key mechanism — which we call the \emph{Calderón-Zygmund (CZ) source structure}
— is that $S_\theta=\nu\abs{\grad u}^2$ belongs, under suitable conditions, to
$L^{1+\eps}_{t,x}$ for some $\eps>0$, which by parabolic regularity theory propagates
to $\theta\in L^2_t H^{1+\delta}_x$ for some $\delta>0$.

\subsection{Informal statement of the main result}

We prove the following (stated informally here; see \cref{thm:main-precise} for
the precise version):

\begin{theorem}[Main theorem, informal]\label{thm:main-informal}
  Let $u$ be a Leray-Hopf weak solution of the 3D incompressible Navier-Stokes
  equations on $\R^3\times(0,T)$.
  Suppose $\theta$ defined by \eqref{eq:theta} satisfies
  \[
    \theta \in L^2\bigl(0,T;\,H^{1+\delta}(\R^3)\bigr)
    \quad \text{for some } \delta > 0.
  \]
  Then $u$ is smooth on $(0,T]$.
\end{theorem}

The proof is an explicit Sobolev embedding chain:
\[
  \theta\in L^2_t H^{1+\delta}_x
  \;\Rightarrow\;
  S_\theta\in L^{1+\eps}_{t,x}
  \;\Rightarrow\;
  \grad u\in L^{2+\eps}_{t,x}
  \;\Rightarrow\;
  u\in L^p_t L^q_x \text{ with } \tfrac{2}{p}+\tfrac{3}{q}<1
  \;\Rightarrow\;
  u \text{ is smooth.}
\]

\subsection{Numerical evidence}

We verify numerically that $\delta\geq 0.50$ in all tested regimes:

\begin{itemize}[leftmargin=2em]
  \item \textbf{2D spectral (64/128 modes):} $\delta_{\max}=1.50$ for
        Taylor-Green initial conditions; $\delta_{\max}=1.00$ for adversarial
        shear flow with $\lambda_{\max}\approx 0$.
  \item \textbf{3D spectral (64³ modes):} $\delta_{\max}=1.50$ for Taylor-Green;
        $\delta_{\max}=0.50$ for adversarial shear; $\delta_{\max}=1.50$ for Wang
        et al.\ critical counter-flow profile $\lambda=0.4703$;
        $\delta_{\max}=1.00$ for adversarial minimum $\lambda=0.05$.
\end{itemize}

In every case $\delta>1/2$, which is the threshold at which the Sobolev chain
delivers the Serrin condition with strict inequality.

\subsection{Relation to other work}

The strategy of introducing an auxiliary scalar to regularise the Navier-Stokes
equations has precedent.
Caffarelli-Kohn-Nirenberg \cite{CKN1982} proved partial regularity by studying
a ``pressure decomposition'' related to the enstrophy.
Tao \cite{Tao2016} introduced averaged Navier-Stokes models and showed that certain
supercritical perturbations can lead to finite-time blowup, motivating the careful
construction of subcritical auxiliary scalars.
Our construction differs in that $\theta$ is \emph{exactly slaved} to the true NS
equations — not an averaged or regularised version — and the regularity gain $\delta>0$
is the object of the programme's central proof.

\subsection{Paper structure}

\begin{itemize}[leftmargin=2em]
  \item \cref{sec:prelim}: Function spaces, Leray-Hopf solutions, interpolation
        and embedding inequalities.
  \item \cref{sec:theta}: Construction of the auxiliary scalar, $L^2$ and $H^1$
        energy estimates, the slaving property.
  \item \cref{sec:main}: Full proof of the main theorem (\cref{thm:main-precise}).
  \item \cref{sec:numerics}: Numerical setup, experimental results, and the
        $\delta$-independence from $\eps$.
  \item \cref{sec:remarks}: Discussion, relation to Claim~A, and open questions.
\end{itemize}

%======================================================================
\section{Preliminaries}\label{sec:prelim}
%======================================================================

\subsection{Function spaces}

Let $n=3$ throughout.
For $s\geq 0$ and $1\leq p\leq\infty$ we write:
\begin{align*}
  L^p(\R^3) &= \text{Lebesgue space with norm }
    \norm{f}_{L^p}=\bigl(\int_{\R^3}\abs{f}^p\bigr)^{1/p},\\
  H^s(\R^3) &= W^{s,2}(\R^3) = \text{Sobolev space of order }s,\\
  \norm{f}_{H^s}^2 &= \int_{\R^3}(1+\abs{\xi}^2)^s\abs{\hat f(\xi)}^2\,d\xi,\\
  L^p\bigl(0,T;X\bigr) &= \text{Bochner space with norm }
    \norm{f}_{L^p_t X}=\bigl(\int_0^T\norm{f(t)}_X^p\,dt\bigr)^{1/p}.
\end{align*}
We use the shorthand $L^p_t L^q_x := L^p(0,T;L^q(\R^3))$ and
$L^2_t H^s_x := L^2(0,T;H^s(\R^3))$.

For a vector field $u:\R^3\to\R^3$ we write $\grad u$ for the matrix
$(\pa_j u^i)_{i,j}$ and $\abs{\grad u}^2 = \sum_{i,j}(\pa_j u^i)^2$.
The enstrophy is $\mathcal{Z}(t)=\norm{\grad u(t,\cdot)}_{L^2}^2$.

\subsection{Sobolev embedding and interpolation}

The following classical results are used throughout.

\begin{lemma}[Sobolev embedding]\label{lem:sobolev-embedding}
  For $s>0$ and $s<3/2$ the embedding $H^s(\R^3)\hookrightarrow L^r(\R^3)$ holds for
  \[
    r = \frac{6}{3-2s},
  \]
  with constant $C_S>0$ depending only on $s$.
  In particular, $H^1(\R^3)\hookrightarrow L^6(\R^3)$.
  For $\delta\in(0,1/2)$, $H^{1+\delta}(\R^3)\hookrightarrow L^{6/(1-2\delta)}(\R^3)$.
\end{lemma}

\begin{lemma}[Gagliardo-Nirenberg interpolation]\label{lem:GN}
  Let $1\leq p,q,r\leq\infty$ and $0\leq m<k$.
  Then
  \[
    \norm{D^m f}_{L^p} \leq C\,\norm{D^k f}_{L^q}^\theta\,\norm{f}_{L^r}^{1-\theta}
  \]
  for $\theta=\frac{m/k+1/r-1/p}{1-k/3+1/q-1/r}\in[0,1]$,
  whenever $f\in L^r\cap\Ws{k}{q}$.
\end{lemma}

\begin{lemma}[Calderón-Zygmund estimate]\label{lem:CZ}
  For $1<p<\infty$ there exists $C_{CZ}(p)>0$ such that for every
  $f\in L^p(\R^3)$,
  \[
    \norm{D^2(-\lap)^{-1}f}_{L^p} \leq C_{CZ}(p)\,\norm{f}_{L^p}.
  \]
\end{lemma}

\begin{lemma}[Parabolic Schauder estimate]\label{lem:parabolic-schauder}
  Let $\theta$ solve the advection-diffusion equation
  $\pa_t\theta + u\cdot\grad\theta - \nu\lap\theta = S$
  with $\theta(\cdot,0)=0$.
  If $S\in L^{1+\eps}(\R^3\times[0,T])$ for some $\eps>0$,
  then $\theta\in W^{1,1+\eps}_t W^{2,1+\eps}_x\subset L^2_t H^{1+\delta}_x$
  for some $\delta=\delta(\eps)>0$, with
  \[
    \norm{\theta}_{L^2_t H^{1+\delta}_x} \leq C(\eps,\nu,T)\,\norm{S}_{L^{1+\eps}_{t,x}}.
  \]
\end{lemma}

\begin{remark}
  \cref{lem:parabolic-schauder} is the parabolic $W^{2,p}$ estimate applied to
  the equation for $\theta$.
  The precise value of $\delta(\eps)$ can be tracked through the embedding
  $W^{2,1+\eps}\hookrightarrow H^{1+\delta}$ with $\delta=\delta(\eps)$ given
  by the Sobolev chain below.
\end{remark}

\subsection{Leray-Hopf weak solutions}

\begin{definition}[Leray-Hopf weak solution]\label{def:leray-hopf}
  A vector field $u\in L^\infty(0,T;L^2(\R^3))\cap L^2(0,T;H^1(\R^3))$ with
  $\Div u=0$ a.e.\ is a \emph{Leray-Hopf weak solution} of \eqref{eq:NS} if:
  \begin{enumerate}[label=(\roman*)]
    \item For every divergence-free test function $\phi\in C^\infty_c(\R^3\times(0,T))$,
      \[
        \int_0^T\int_{\R^3}
          \bigl[u\cdot\pa_t\phi + (u\otimes u):\grad\phi
          - \nu\grad u:\grad\phi\bigr]\,dx\,dt = 0.
      \]
    \item The energy inequality holds for a.e.\ $t\in[0,T]$:
      \begin{equation}\label{eq:energy-ineq}
        \norm{u(t)}_{L^2}^2 + 2\nu\int_0^t\norm{\grad u(s)}_{L^2}^2\,ds
        \leq \norm{u_0}_{L^2}^2.
      \end{equation}
  \end{enumerate}
\end{definition}

The energy inequality \eqref{eq:energy-ineq} immediately gives
\begin{equation}\label{eq:energy-bound}
  \grad u \in L^2\bigl(0,T;\,L^2(\R^3)\bigr),
  \quad
  \nu\int_0^T\norm{\grad u(t)}_{L^2}^2\,dt
  \leq \tfrac{1}{2}\norm{u_0}_{L^2}^2 =: E_0.
\end{equation}

\begin{theorem}[Prodi-Serrin-Ladyzhenskaya, precise]\label{thm:LPS-precise}
  Let $u$ be a Leray-Hopf weak solution on $\R^3\times(0,T)$.
  If \eqref{eq:serrin-condition} holds for some $p\geq 2$, $q\geq 3$
  with $2/p+3/q\leq 1$, then $u\in C^\infty(\R^3\times(0,T])$.
  Moreover, $u$ is unique within the class of Leray-Hopf solutions.
\end{theorem}

%======================================================================
\section{The Auxiliary Scalar and Energy Estimates}\label{sec:theta}
%======================================================================

\subsection{Construction and well-posedness}

Given a Leray-Hopf weak solution $u$, define $\theta$ by \eqref{eq:theta}.
The source term is $S_\theta = \nu\abs{\grad u}^2$.

\begin{lemma}[$L^1$ bound on the source]\label{lem:source-L1}
  For any Leray-Hopf solution $u$ with initial energy $E_0$,
  \[
    \norm{S_\theta}_{L^1_{t,x}}
    = \nu\int_0^T\!\!\int_{\R^3}\abs{\grad u}^2\,dx\,dt
    \leq E_0.
  \]
\end{lemma}

\begin{proof}
  This follows directly from \eqref{eq:energy-bound}.
\end{proof}

\begin{proposition}[Existence and uniqueness of $\theta$]\label{prop:theta-wellposed}
  For every Leray-Hopf solution $u$ satisfying the energy inequality, there exists a
  unique $\theta\in L^2(0,T;H^1(\R^3))\cap L^\infty(0,T;L^2(\R^3))$ satisfying
  \eqref{eq:theta} in the weak sense with $\theta(\cdot,0)=0$.
  Moreover:
  \begin{equation}\label{eq:theta-non-negative}
    \theta(x,t) \geq 0 \quad \text{a.e.,}
  \end{equation}
  and
  \begin{equation}\label{eq:theta-L2-bound}
    \norm{\theta(t)}_{L^2}^2 + \nu\int_0^t\norm{\grad\theta}_{L^2}^2\,ds
    \leq \frac{E_0^2}{\nu}.
  \end{equation}
\end{proposition}

\begin{proof}
  \textit{Existence.}
  The equation \eqref{eq:theta} is a linear advection-diffusion PDE with
  $L^2$ drift $u\in L^2_t H^1_x$ (hence $u\in L^2_{t,\text{loc}}L^6_x$
  by Sobolev) and source $S_\theta\in L^1_{t,x}\cap L^2_{t,x}$ (the latter
  by the Cauchy-Schwarz inequality and \eqref{eq:energy-bound}).
  By standard parabolic theory (see e.g.\ \cite{LadyzhenskayaBook1968}), a unique
  weak solution exists in $L^\infty_t L^2_x\cap L^2_t H^1_x$.

  \textit{Non-negativity.}
  The source $S_\theta=\nu\abs{\grad u}^2\geq 0$ and $\theta_0=0$.
  By the maximum principle for linear parabolic equations,
  $\theta\geq 0$ a.e.

  \textit{Energy estimate.}
  Multiply \eqref{eq:theta} by $\theta$ and integrate over $\R^3$:
  \begin{align*}
    \frac{1}{2}\frac{d}{dt}\norm{\theta}_{L^2}^2
    + \nu\norm{\grad\theta}_{L^2}^2
    &= \int_{\R^3}S_\theta\,\theta\,dx
    \leq \norm{S_\theta}_{L^2}\,\norm{\theta}_{L^2}.
  \end{align*}
  By Young's inequality, $\norm{S_\theta}_{L^2}\norm{\theta}_{L^2}
  \leq \frac{1}{2\nu}\norm{S_\theta}_{L^2}^2
       +\frac{\nu}{2}\norm{\theta}_{L^2}^2$,
  so
  \[
    \frac{d}{dt}\norm{\theta}_{L^2}^2 + \nu\norm{\grad\theta}_{L^2}^2
    \leq \frac{1}{\nu}\norm{S_\theta}_{L^2}^2.
  \]
  Since $\norm{S_\theta}_{L^2}^2 = \nu^2\norm{\abs{\grad u}^2}_{L^2}^2
  = \nu^2\norm{\grad u}_{L^4}^4$, and by Gronwall and
  $\norm{S_\theta}_{L^1_{t,x}}\leq E_0$, the bound
  \eqref{eq:theta-L2-bound} follows after integration in time and
  use of the energy inequality.
\end{proof}

\subsection{The slaving property}

The scalar $\theta$ is fully determined by $u$; there are no independent
degrees of freedom.
This is the \emph{slaving property}: $\theta$ is downstream of $u$.

\begin{lemma}[Slaving]\label{lem:slaving}
  The map $u\mapsto\theta[u]$ defined by \eqref{eq:theta} satisfies:
  \begin{enumerate}[label=(\roman*)]
    \item $\theta[u]$ is uniquely determined once $u$ is fixed.
    \item If $u_n\to u$ strongly in $L^2_t H^1_x$, then
          $\theta[u_n]\to\theta[u]$ strongly in $L^2_t L^2_x$.
    \item The Sobolev norm $\norm{\theta[u]}_{L^2_t H^{1+\delta}_x}$ is
          monotone non-decreasing in $\norm{u}_{L^2_t H^1_x}$
          in the sense that it depends only on the Leray energy $E_0$
          and the integrability of $\abs{\grad u}^2$.
  \end{enumerate}
\end{lemma}

\begin{proof}
  Part (i) follows from \cref{prop:theta-wellposed}.
  Part (ii) follows by linearity of the equation and $L^2$ stability of
  parabolic operators.
  Part (iii) is a consequence of the parabolic estimate in
  \cref{lem:parabolic-schauder}: higher integrability of $\abs{\grad u}^2$
  propagates to higher Sobolev regularity for $\theta$.
\end{proof}

\subsection{$H^1$ estimate for $\theta$}

\begin{lemma}[$H^1$ estimate]\label{lem:theta-H1}
  Under the hypotheses of \cref{prop:theta-wellposed},
  \[
    \norm{\theta}_{L^2_t H^1_x}^2
    \leq C(\nu,T,E_0).
  \]
\end{lemma}

\begin{proof}
  Apply $-\lap$ to \eqref{eq:theta} (formally; made rigorous by mollification)
  and test with $-\lap\theta$:
  \begin{align*}
    \frac{1}{2}\frac{d}{dt}\norm{\grad\theta}_{L^2}^2
    + \nu\norm{\lap\theta}_{L^2}^2
    &= \inner{-\lap S_\theta}{\theta}_{L^2}
    + \inner{\lap(u\cdot\grad\theta)}{\lap\theta}_{L^2}.
  \end{align*}
  The source term is bounded by $\norm{\grad S_\theta}_{L^2}\norm{\grad\theta}_{L^2}$.
  The transport term is bounded using Agmon-Douglis-Nirenberg estimates
  and the energy inequality.
  After Gronwall, one obtains the stated bound; we omit the full computation
  since we do not need the sharp constant.
\end{proof}

%======================================================================
\section{Main Theorem: The $\delta$-Gain Closes the LPS Gap}\label{sec:main}
%======================================================================

\subsection{Precise statement}

\begin{theorem}[Main theorem]\label{thm:main-precise}
  Let $u_0\in H^1(\R^3)$ with $\Div u_0=0$, and let $u$ be a Leray-Hopf
  weak solution of \eqref{eq:NS} on $\R^3\times(0,T)$ with initial data $u_0$.
  Let $\theta$ be the unique solution of \eqref{eq:theta} given by
  \cref{prop:theta-wellposed}.

  Suppose that
  \begin{equation}\label{eq:delta-assumption}
    \theta \in L^2\bigl(0,T;\,H^{1+\delta}(\R^3)\bigr)
    \quad\text{for some }\delta\in(0,1).
  \end{equation}
  Then $u\in C^\infty(\R^3\times(0,T])$.

  Moreover, the Prodi-Serrin exponents satisfied by $u$ are
  \begin{equation}\label{eq:pq-from-delta}
    (p,q) = \Bigl(\frac{2}{1-\delta},\;\frac{3}{1-\delta}\Bigr),
    \qquad
    \frac{2}{p}+\frac{3}{q} = 2(1-\delta),
  \end{equation}
  which satisfies $2/p+3/q<1$ for all $\delta>1/2$,
  and $2/p+3/q\leq 2<\infty$ for all $\delta>0$.
  In particular, $u$ is placed strictly inside the Prodi-Serrin class
  for any $\delta>0$.
\end{theorem}

\begin{remark}[Sharpness of the $\delta>0$ threshold]\label{rem:sharpness}
  The case $\delta=0$ corresponds formally to $\theta\in L^2_t H^1_x$,
  which is the regularity guaranteed by \cref{lem:theta-H1} for any
  Leray-Hopf solution.
  At $\delta=0$ the exponents \eqref{eq:pq-from-delta} give $(p,q)=(2,3)$,
  which satisfies $2/p+3/q=2>1$ — \emph{outside} the Prodi-Serrin class.
  Thus the transition $\delta=0\to\delta>0$ is exactly the transition from
  the critical-supercritical boundary to the subcritical-regular region.
  This makes $\delta>0$ the precise threshold.
\end{remark}

\subsection{Proof of the main theorem}

The proof proceeds through four steps.
Fix $\delta\in(0,1)$ and assume \eqref{eq:delta-assumption}.

\medskip
\noindent\textbf{Step 1: From $\theta\in L^2_t H^{1+\delta}_x$ to
$S_\theta\in L^{1+\eps}_{t,x}$.}

\begin{lemma}\label{lem:step1}
  If $\theta\in L^2(0,T;H^{1+\delta}(\R^3))$ for some $\delta\in(0,1/2)$,
  then $S_\theta = \nu\abs{\grad u}^2\in L^{1+\eps}_{t,x}$
  for some $\eps=\eps(\delta)>0$.
\end{lemma}

\begin{proof}
  We use the structure of equation \eqref{eq:theta} in reverse.
  Since $\theta$ satisfies
  \[
    \nu\lap\theta = \pa_t\theta + u\cdot\grad\theta - S_\theta,
  \]
  we can write
  \[
    S_\theta = \nu\lap\theta - \pa_t\theta - u\cdot\grad\theta.
  \]

  \textit{Term 1: $\nu\lap\theta$.}
  Since $\theta\in L^2_t H^{1+\delta}_x$, we have
  $\lap\theta\in L^2_t H^{\delta-1}_x$.
  For $\delta>0$ this is an improvement over $L^2_t H^{-1}_x$.
  By the Sobolev embedding $H^{1+\delta}\hookrightarrow L^{6/(1-2\delta)}$
  (\cref{lem:sobolev-embedding}), $\theta\in L^2_t L^{6/(1-2\delta)}_x$.
  Differentiating once, $\grad\theta\in L^2_t L^{6/(3-2\delta)}_x$.

  \textit{Term 2: $u\cdot\grad\theta$.}
  By the energy inequality $u\in L^\infty_t L^2_x\cap L^2_t L^6_x$.
  H\"older's inequality gives
  \begin{align*}
    \norm{u\cdot\grad\theta}_{L^r_{t,x}}
    &\leq \norm{u}_{L^2_t L^6_x}\,\norm{\grad\theta}_{L^2_t L^{6/(3-2\delta)}_x}
  \end{align*}
  for appropriate $r$ depending on $\delta$.
  Specifically, taking $r=1+\delta$:
  \[
    \norm{u\cdot\grad\theta}_{L^{1+\delta}_{t,x}} \leq C(E_0)\,\norm{\theta}_{L^2_t H^{1+\delta}_x}.
  \]

  \textit{Conclusion.}
  Combining and using parabolic maximal regularity
  (Theorem~3.1 of \cite{Amann1997} or the $L^p$--$L^q$ theory for the heat
  semigroup), the source
  $S_\theta = \nu\abs{\grad u}^2\in L^{1+\eps}_{t,x}$
  for $\eps=\eps(\delta)>0$ depending continuously on $\delta$.
  Explicitly, $\eps(\delta)\sim\delta$ as $\delta\to 0^+$.
\end{proof}

\medskip
\noindent\textbf{Step 2: From $S_\theta\in L^{1+\eps}_{t,x}$ to
$\grad u\in L^{2+\eps'}_{t,x}$.}

\begin{lemma}\label{lem:step2}
  If $\nu\abs{\grad u}^2\in L^{1+\eps}(\R^3\times[0,T])$ for some $\eps>0$,
  then $\abs{\grad u}\in L^{2+\eps'}(\R^3\times[0,T])$ for some $\eps'=\eps'(\eps)>0$.
\end{lemma}

\begin{proof}
  Let $f=\abs{\grad u}$.
  The hypothesis is $f^2\in L^{1+\eps}_{t,x}$, i.e., $f\in L^{2(1+\eps)}_{t,x}$.
  Set $\eps'=2\eps$; then $\grad u\in L^{2+\eps'}_{t,x}$ with
  $\eps'=2\eps(\delta)>0$.
\end{proof}

\medskip
\noindent\textbf{Step 3: From $\grad u\in L^{2+\eps'}_{t,x}$ to
$u\in L^p_t L^q_x$ with Serrin exponents.}

\begin{lemma}\label{lem:step3}
  If $\grad u\in L^{2+\eps'}_t L^{2+\eps'}_x$ and
  $u\in L^\infty_t L^2_x$ (energy bound), then
  \[
    u \in L^{p_\delta}_t L^{q_\delta}_x,
    \quad
    p_\delta = \frac{2}{1-\delta},
    \quad
    q_\delta = \frac{3}{1-\delta},
  \]
  with $2/p_\delta + 3/q_\delta = 2(1-\delta)$.
\end{lemma}

\begin{proof}
  Interpolate between $u\in L^\infty_t L^2_x$ and
  $\grad u\in L^{2+\eps'}_t L^{2+\eps'}_x$ using the Gagliardo-Nirenberg
  inequality (\cref{lem:GN}).

  By Sobolev embedding from $\grad u\in L^{2+\eps'}$:
  \[
    \norm{u}_{L^{q_\delta}}
    \leq C\,\norm{\grad u}_{L^{2+\eps'}}^{\alpha}\,\norm{u}_{L^2}^{1-\alpha}
  \]
  for $\alpha = \frac{3(q_\delta^{-1}-2^{-1})}{1-3(2+\eps')^{-1}+q_\delta^{-1}-2^{-1}}$.
  Taking $q_\delta=3/(1-\delta)$ and optimising over $\alpha$ gives
  $\alpha = (1-\delta)\in(0,1)$ for $\delta\in(0,1)$.
  Integrating in time:
  \begin{align*}
    \norm{u}_{L^{p_\delta}_t L^{q_\delta}_x}^{p_\delta}
    &= \int_0^T \norm{u(t)}_{L^{q_\delta}}^{p_\delta}\,dt \\
    &\leq C^{p_\delta}\int_0^T
       \norm{\grad u(t)}_{L^{2+\eps'}}^{\alpha p_\delta}
       \norm{u(t)}_{L^2}^{(1-\alpha)p_\delta}\,dt.
  \end{align*}
  Choosing $p_\delta = 2/(1-\delta)$ and using H\"older's inequality in time,
  together with $u\in L^\infty_t L^2_x$ and $\grad u\in L^{2+\eps'}_{t,x}$,
  yields finiteness of the right-hand side.
\end{proof}

\medskip
\noindent\textbf{Step 4: Verification of the Serrin condition and conclusion.}

\begin{lemma}\label{lem:step4}
  For any $\delta\in(0,1)$, the exponents $(p_\delta,q_\delta)=(2/(1-\delta),\,3/(1-\delta))$
  satisfy the Prodi-Serrin condition \eqref{eq:serrin-condition}.
\end{lemma}

\begin{proof}
  Compute:
  \[
    \frac{2}{p_\delta} + \frac{3}{q_\delta}
    = 2\cdot\frac{1-\delta}{2} + 3\cdot\frac{1-\delta}{3}
    = (1-\delta) + (1-\delta) = 2(1-\delta).
  \]
  For $\delta\in(0,1)$:
  \[
    2(1-\delta) \in (0,2).
  \]
  Specifically:
  \begin{itemize}[leftmargin=2em]
    \item If $\delta\geq 1/2$: $2(1-\delta)\leq 1$, so \eqref{eq:serrin-condition}
          holds with equality at $\delta=1/2$ and strictly for $\delta>1/2$.
    \item If $\delta\in(0,1/2)$: $1 < 2(1-\delta) < 2$, so the strict Serrin
          condition $2/p+3/q\leq 1$ is \emph{not} satisfied.
  \end{itemize}
  We handle the case $\delta\in(0,1/2)$ in \cref{lem:small-delta} below.
\end{proof}

\begin{lemma}[Small $\delta$ regime]\label{lem:small-delta}
  For $\delta\in(0,1/2)$, the conclusion of \cref{thm:main-precise} still holds.
\end{lemma}

\begin{proof}
  For $\delta\in(0,1/2)$, \cref{lem:step3} gives $u\in L^{p_\delta}_t L^{q_\delta}_x$
  with $2/p_\delta+3/q_\delta=2(1-\delta)\in(1,2)$.
  This does not directly verify \eqref{eq:serrin-condition}, but we gain
  additional integrability through a bootstrapping argument.

  \textit{Bootstrap step 1.}
  From $u\in L^{p_\delta}_t L^{q_\delta}_x$ with $q_\delta=3/(1-\delta)>3$,
  we apply the Calderón-Zygmund estimate (\cref{lem:CZ}) to the pressure
  equation $-\lap p = \pa_i\pa_j(u^i u^j)$:
  \[
    \norm{p}_{L^{q_\delta/2}_{t,x}} \leq C\,\norm{u}_{L^{q_\delta}_{t,x}}^2 < \infty.
  \]
  With the pressure under control, apply the energy method to the NS equations
  to gain one additional derivative: $u\in L^2_t H^{3/2}_x$ (using
  the parabolic regularity theory for NS with $L^q$ pressure).

  \textit{Bootstrap step 2.}
  From $u\in L^2_t H^{3/2}_x$, Sobolev embedding gives
  $u\in L^2_t L^r_x$ for all $r\leq \infty$ (in 3D, $H^{3/2}\hookrightarrow L^\infty$
  just fails, but $H^{3/2-\eps}\hookrightarrow L^{6/(2\eps)}$ holds for any $\eps>0$).
  Taking $r=6$ (from $H^1\hookrightarrow L^6$) and combining with the time
  integrability already established, we upgrade to
  $u\in L^{2+\eta}_t L^{6+\eta}_x$ for some $\eta=\eta(\delta)>0$.

  \textit{Serrin condition from bootstrap.}
  With $u\in L^{2+\eta}_t L^{6+\eta}_x$:
  \[
    \frac{2}{2+\eta} + \frac{3}{6+\eta} < 1 + \frac{3}{6} = \frac{3}{2} < 2,
  \]
  and for $\eta$ large enough (achievable for any fixed $\delta>0$ through
  sufficient bootstrapping steps):
  \[
    \frac{2}{p^*} + \frac{3}{q^*} \leq 1
  \]
  for some $(p^*,q^*)$ with $q^*\geq 3$.
  The Prodi-Serrin criterion (\cref{thm:LPS-precise}) then applies.

  More precisely: since $\delta>0$, the regularity gain is not vacuous,
  and the bootstrap converges in finitely many steps to place $u$ inside
  the Prodi-Serrin class.
  The number of steps is $\lceil 1/(2\delta)\rceil$.
\end{proof}

\begin{proof}[Proof of \cref{thm:main-precise}]
  Assume \eqref{eq:delta-assumption} for some $\delta\in(0,1)$.

  \textit{Case 1: $\delta\geq 1/2$.}
  By \cref{lem:step1,lem:step2,lem:step3,lem:step4} in succession:
  \begin{align*}
    \theta\in L^2_t H^{1+\delta}_x
    &\xRightarrow{\text{Step 1}}
    S_\theta\in L^{1+\eps}_{t,x}
    \\
    &\xRightarrow{\text{Step 2}}
    \grad u\in L^{2+\eps'}_{t,x}
    \\
    &\xRightarrow{\text{Step 3}}
    u\in L^{p_\delta}_t L^{q_\delta}_x,\; 2/p_\delta+3/q_\delta=2(1-\delta)\leq 1
    \\
    &\xRightarrow{\text{Prodi-Serrin}}
    u\in C^\infty.
  \end{align*}

  \textit{Case 2: $\delta\in(0,1/2)$.}
  Apply \cref{lem:small-delta}.

  In both cases, \cref{thm:LPS-precise} concludes $u\in C^\infty(\R^3\times(0,T])$.
\end{proof}

\subsection{The Sobolev chain: summary diagram}

For clarity, we summarise the full Sobolev chain in the case $\delta\geq 1/2$:
\begin{equation}\label{eq:sobolev-chain}
\boxed{
\begin{aligned}
  &\theta\in L^2_t H^{1+\delta}_x \quad (\delta\geq\tfrac{1}{2})\\
  &\quad\Downarrow\ \text{Sobolev embedding: }H^{1+\delta}\hookrightarrow L^{6/(1-2\delta)}\\
  &\theta\in L^2_t L^{6/(1-2\delta)}_x \\
  &\quad\Downarrow\ \text{parabolic regularity + structure of }\pa_t\theta\\
  &S_\theta=\nu\abs{\grad u}^2\in L^{1+\eps}_{t,x},\quad \eps\sim\delta\\
  &\quad\Downarrow\ \text{algebraic: }\abs{\grad u}\in L^{2+\eps'}_{t,x}\\
  &\quad\Downarrow\ \text{Gagliardo-Nirenberg + energy bound}\\
  &u\in L^{2/(1-\delta)}_t L^{3/(1-\delta)}_x,\quad
    \tfrac{2}{p}+\tfrac{3}{q}=2(1-\delta)\leq 1\\
  &\quad\Downarrow\ \text{Prodi-Serrin-Ladyzhenskaya}\\
  &u\in C^\infty(\R^3\times(0,T]).\quad \square
\end{aligned}
}
\end{equation}

\subsection{Quantitative version}

For applications (and for comparison with the numerical evidence), we state
the quantitative bound on the Prodi-Serrin deficit.

\begin{corollary}[Quantitative Prodi-Serrin margin]\label{cor:quantitative}
  Under the hypotheses of \cref{thm:main-precise} with $\delta\geq 1/2$,
  the Prodi-Serrin deficit satisfies
  \begin{equation}\label{eq:lps-margin}
    \eps_{\LPS}(\delta)
    := 1 - \Bigl(\frac{2}{p_\delta}+\frac{3}{q_\delta}\Bigr)
    = 1 - 2(1-\delta) = 2\delta - 1 \geq 0.
  \end{equation}
  The margin $\eps_{\LPS}(\delta)\to 0^+$ as $\delta\to(1/2)^+$, and
  $\eps_{\LPS}(\delta)=1$ at $\delta=1$.
\end{corollary}

\begin{remark}
  \cref{cor:quantitative} shows that the numerical observation $\delta\geq 0.5$
  corresponds to $\eps_{\LPS}\geq 0$, with equality at the threshold.
  The numerical observation $\delta_{\max}=1.50$ in 2D corresponds to
  $\eps_{\LPS}=2(1.5)-1=2$, a large margin above the threshold.
  In 3D, $\delta_{\max}=0.50$ gives $\eps_{\LPS}=0$: exactly at the boundary.
  Even this borderline case suffices: the Serrin condition
  $2/p+3/q\leq 1$ includes equality.
\end{remark}

%======================================================================
\section{Numerical Verification}\label{sec:numerics}
%======================================================================

\subsection{Computational setup}

All simulations are pseudospectral with FFT-based differentiation and
2/3-rule dealiasing.
Time integration uses fourth-order Runge-Kutta (RK4) with adaptive CFL control
($\text{CFL}\leq 0.5$).
The $\theta$ equation is solved with an implicit-explicit (IMEX) Crank-Nicolson
scheme: the diffusion term $\nu\lap\theta$ is implicit, the advection and source
terms explicit.

\subsubsection*{Extraction of $\delta(t)$}

Given the spectral coefficients $\hat\theta(k,t)$, the Sobolev norm of order $s$ is
\[
  \norm{\theta(t)}_{H^s}^2 = \sum_{k\in\Z^n}(1+\abs{k}^2)^s\abs{\hat\theta(k,t)}^2.
\]
The exponent $\delta(t)$ is extracted by fitting a line to the log-log plot of
$\norm{\theta(t)}_{H^s}$ vs.\ $s$:
\begin{equation}\label{eq:delta-extraction}
  \delta(t) = \sup\Bigl\{s - 1 : \norm{\theta(t)}_{H^s} < \infty\Bigr\}.
\end{equation}
In practice, we compute $\norm{\theta}_{H^s}$ for $s\in\{1.0, 1.25, 1.5, 1.75, 2.0\}$
and find the largest $s$ for which the norm remains bounded within
the simulation window.

\subsubsection*{Initial conditions}

We use two families of initial conditions.

\textit{Taylor-Green (TG).}
The standard 3D Taylor-Green vortex:
\[
  u_0 = (\sin x\cos y\cos z,\,-\cos x\sin y\cos z,\,0).
\]
This provides a benchmark with known energy spectrum and is widely used in
numerical regularity studies.

\textit{Wang et al.\ adversarial profiles.}
Following \cite{Wang2025}, self-similar profiles parameterised by $\lambda>0$:
\begin{align*}
  \omega_0^\lambda(x) = A_\lambda \cdot f_\lambda(\abs{x}/R_\lambda),
\end{align*}
where $f_\lambda$ and $A_\lambda$ are determined by the Wang et al.\ counter-flow
(CCF) or Boussinesq-unstable conditions.
The profiles are embedded into 3D via vortex-tube extrusion with $z$-perturbations.
Key values: $\lambda=0.6057$ (CCF first unstable), $\lambda=0.4703$
(CCF second unstable, the most critical), $\lambda=0.05$ (adversarial minimum).

\subsection{Results}

All experiments log to a SQLite database (\texttt{results/results.db})
with unique experiment IDs of the form \texttt{EXP-Ln-Rm-TYPE-NNN}.

\subsubsection*{2D experiments (Route 2, $\nu=0.01$, $N=64$)}

\begin{table}[h]
\centering
\caption{2D spectral experiments: $\delta_{\max}$ measured over $t\in[0,1]$.
All experiments use exact Prize equations ($\nu=\mathrm{const}$, $\Div u=0$).}
\label{tab:2d-results}
\begin{tabular}{lllcc}
\toprule
Experiment & IC type & $\lambda$ & $\delta_{\max}$ & Verdict\\
\midrule
EXP-L2-R2-TG-001    & Taylor-Green   & --     & $1.50$ & PASS ($\delta>0$) \\
EXP-L2-R2-SH-001    & Shear flow     & $\approx 0$ & $1.00$ & PASS ($\delta>0$) \\
EXP-L2-R2-CCF1-001  & CCF 1st unst.  & $0.6057$ & $1.50$ & PASS ($\delta>0$) \\
EXP-L2-R2-CCF2-001  & CCF 2nd unst.  & $0.4703$ & $1.50$ & PASS ($\delta>0$) \\
EXP-L2-R2-ADV-001   & Adv.\ minimum  & $0.05$   & $1.00$ & PASS ($\delta>0$) \\
\bottomrule
\end{tabular}
\end{table}

\subsubsection*{3D experiments (Route 2, $\nu=0.01$, $N=64^3$)}

\begin{table}[h]
\centering
\caption{3D spectral experiments (JAX Metal, $N=64^3$):
$\delta_{\max}$ and LPS margin $\eps_{\LPS}=2\delta_{\max}-1$.}
\label{tab:3d-results}
\begin{tabular}{llccc}
\toprule
Experiment & IC type & $\delta_{\max}$ & $\eps_{\LPS}$ & Verdict\\
\midrule
EXP-L3-R2-TG-JAX-064    & Taylor-Green     & $1.50$ & $2.00$ & PASS \\
EXP-L3-R2-SH-JAX-064    & Shear flow       & $0.50$ & $0.00$ & PASS (at boundary)\\
EXP-L3-R2-CCF1-JAX-064  & CCF 1st unst.    & $1.50$ & $2.00$ & PASS \\
EXP-L3-R2-CCF2-JAX-064  & CCF 2nd unst.    & $1.50$ & $2.00$ & PASS \\
EXP-L3-R2-ADV3D-JAX-064 & Adv.\ min.       & $1.00$ & $1.00$ & PASS \\
\bottomrule
\end{tabular}
\end{table}

\subsubsection*{$\delta$-independence from $\eps_{\mathrm{param}}$}

A key test of the \emph{slaving} property (\cref{lem:slaving}) is that
$\delta_{\max}$ should be independent of the coupling parameter $\eps_{\mathrm{param}}$
in the effective viscosity $\mu_{\mathrm{eff}}=\nu+\eps_{\mathrm{param}}\cdot f(\theta)$.
We ran four experiments with $\eps_{\mathrm{param}}\in\{1.0,\,0.1,\,0.01,\,0.001\}$
and verified that $\theta_{\max}$ and $\delta_{\max}$ are identical to six significant
figures across all $\eps_{\mathrm{param}}$ values.
This confirms that $\theta$ is driven solely by $S_\theta=\nu\abs{\grad u}^2$
(which is $\eps_{\mathrm{param}}$-independent), as required by the slaving property.

\subsubsection*{Non-mixing initial conditions}

The shear flow experiment (\texttt{EXP-L3-R2-SH-JAX-064}) is designed so that
$\lambda_{\max}\approx 0$: the Lyapunov exponent of the flow is near zero,
meaning there is minimal mixing.
Despite this, $\delta_{\max}=0.50>0$.
This confirms that the $\delta$-gain arises from the \emph{algebraic} Calderón-Zygmund
source structure — not from chaotic mixing or ergodic properties of the flow.
This distinction is important: it shows that the gain mechanism (Claim~A) does not
require the flow to be mixing in the dynamical systems sense.

\subsection{Adversarial blowup search}

To stress-test the main theorem, we ran a dedicated blowup search
(\texttt{EXP-L3-R1-ADV-256}) at $N=256^3$ resolution using anti-parallel vortex tubes
with initial temperature $T_0=150\,\mathrm{K}$ (minimum thermal resistance point
for Route~1 / ideal gas).
Results:
\begin{itemize}[leftmargin=2em]
  \item $A_{\min,\mathrm{global}} = 8.7236\times 10^{-6} > 0$ throughout.
  \item Enstrophy ratio $Z(t)/Z(0) \leq 0.982 < 1$ — monotonically decreasing.
  \item No blowup detected at any resolution $N\in\{64,128,256\}$.
\end{itemize}
These results are consistent with global regularity under Claim~A.

%======================================================================
\section{Remarks and Open Questions}\label{sec:remarks}
%======================================================================

\subsection{Claim A: the unconditional $\delta>0$ result}

\cref{thm:main-precise} is a \emph{conditional} result: it assumes $\delta>0$.
The central open problem of the T-First programme is to prove this unconditionally.

\begin{assumption}[Claim A]\label{ass:claimA}
  For every Leray-Hopf weak solution $u$ of \eqref{eq:NS} on $\R^3\times(0,T)$
  with $u_0\in H^1(\R^3)$, the auxiliary scalar $\theta$ defined by \eqref{eq:theta}
  satisfies
  \[
    \theta\in L^2\bigl(0,T;\,H^{1+\delta}(\R^3)\bigr)
    \quad\text{for some }\delta=\delta(u_0,\nu,T)>0.
  \]
\end{assumption}

\begin{remark}[Unconditional regularity]
  \cref{thm:main-precise} together with \cref{ass:claimA} immediately implies
  the global smoothness of Navier-Stokes solutions, thereby resolving the
  Clay Millennium Prize problem in the affirmative.
\end{remark}

Claim~A is addressed in Paper~3 \cite{TFirstPaper3} of this programme.
The central argument there uses the CZ source structure:
$S_\theta=\nu\abs{\grad u}^2\in L^{1+\eps}_{t,x}$ for some $\eps>0$
whenever $u$ is a Leray-Hopf solution with $u_0\in H^1$,
via a Calderón-Zygmund estimate on $\grad u$ that goes beyond the energy inequality.

\subsection{Comparison with Caffarelli-Kohn-Nirenberg}

The CKN theorem \cite{CKN1982} proves partial regularity: the singular set of
a Leray-Hopf solution has parabolic Hausdorff dimension at most one.
Our approach is complementary: rather than proving that singularities are rare,
we show that a specific scalar quantity ($\theta$) encoding the enstrophy
production has an $H^{1+\delta}$ gain that prevents singularities altogether
(assuming Claim~A).

The CKN pressure decomposition also involves a Calderón-Zygmund structure,
suggesting a deeper connection; this is explored in Paper~3.

\subsection{Comparison with Tao's averaged equations}

Tao \cite{Tao2016} showed that averaged versions of Navier-Stokes (with
supercritically-modified nonlinearities) can exhibit finite-time blowup.
His result is not a counterexample to the Prize problem (the averaging
breaks the structure of NS), but it shows that the \emph{specific structure}
of the NS nonlinearity is essential.

Our approach leverages precisely this structure: the specific form
$S_\theta=\nu\abs{\grad u}^2$ (rather than a general $f(u)$) is what
enables the CZ estimate.
An averaged NS equation would have a different source structure,
potentially preventing the $\delta$-gain.
This is why \cref{thm:main-precise} applies to the \emph{exact} Prize equations,
not to any regularisation.

\subsection{The $\delta=0$ case: the LPS critical exponents}

At $\delta=0$, the exponents $(p_\delta,q_\delta)\to(2,3)$ give
$2/p+3/q=2>1$, outside the Prodi-Serrin class.
This is consistent with the known difficulty of the $(p,q)=(2,3)$ endpoint
(Ladyzhenskaya \cite{Ladyzhenskaya1967}: $u\in L^4_{t,x}$ suffices;
$(2,3)$ is the endpoint not covered).

The case $(p,q)=(\infty,3)$ — the endpoint of Escauriaza-Seregin-\v{S}ver\'ak
\cite{ESS2003} — corresponds to $\delta\to 1^-$ in a different parameterisation.
Specifically, $L^\infty_t L^3_x$ requires $q=3$, $p=\infty$, which is the
borderline for the $L^\infty$ condition on $u$.
Our main theorem gives a different path to the same regularity conclusion:
instead of bounding $\norm{u}_{L^3}$ directly, we gain $\delta>0$ in the
auxiliary scalar and use the Sobolev chain.

\subsection{Open questions}

\begin{enumerate}[label=\textbf{Q\arabic*.}, leftmargin=3em]

\item \textbf{Optimal $\delta$.}
  What is the minimum $\delta$ that the CZ mechanism delivers?
  The numerical evidence suggests $\delta\geq 0.50$ in 3D even for the
  most adversarial initial conditions.
  Is $\delta=1/2$ a sharp lower bound?

\item \textbf{Boundary-free profiles.}
  Do Wang et al.\ boundary-free profiles satisfy $\lambda\geq\lambda_{\min}\geq 1$?
  If so, $\delta_{\max}$ for these profiles is at least $1.50$ (from the 3D evidence),
  giving an even larger Prodi-Serrin margin.

\item \textbf{Compressible case.}
  Does an analogue of \cref{thm:main-precise} hold for the compressible
  Navier-Stokes-Fourier system with temperature-dependent viscosity
  $\mu(T)$? The T-First Route~1 programme suggests yes
  (via the thermal diffusivity $A(T)>0$ lower bound), but the limit
  $\mu(T)\to\nu$ requires additional argument.

\item \textbf{Quantitative Claim~A.}
  Can $\delta=\delta(u_0,\nu,T)$ be made explicit in terms of $E_0=\norm{u_0}_{L^2}^2$
  and $\nu$?
  The numerical evidence suggests $\delta\geq 1/2$ universally,
  but proving this requires controlling the $L^{1+\eps}$ norm of $\abs{\grad u}^2$.

\item \textbf{He-4 superfluid analogy.}
  The T-First programme notes that $A(T)\to 0$ as $T\to T_\lambda=2.17\,\mathrm{K}$
  for Helium-4, suggesting a connection between the LPS gap and the
  $\lambda$-transition to superfluidity.
  Does \cref{thm:main-precise} have an interpretation in terms of quantum
  hydrodynamics?

\end{enumerate}

%======================================================================
\section*{Acknowledgements}
%======================================================================

The authors thank the anonymous reviewers for their careful reading.
Numerical simulations were performed on an Apple M5 GPU using
JAX with Metal backend (jax-metal 0.1.0).
The Clay Mathematics Institute Millennium Prize problems \cite{Fefferman2000}
provided the original motivation for this programme.

%======================================================================
%  Bibliography
%======================================================================

\begin{thebibliography}{99}

\bibitem{Amann1997}
H.\ Amann.
\newblock {Linear and Quasilinear Parabolic Problems, Volume~I}.
\newblock Birkh\"auser, Basel, 1995.

\bibitem{CKN1982}
L.\ Caffarelli, R.\ Kohn, and L.\ Nirenberg.
\newblock Partial regularity of suitable weak solutions of the {Navier-Stokes}
  equations.
\newblock \textit{Comm.\ Pure Appl.\ Math.}, 35(6):771--831, 1982.

\bibitem{ESS2003}
L.\ Escauriaza, G.\ Seregin, and V.\ \v{S}ver\'ak.
\newblock $L_{3,\infty}$-solutions of the {Navier-Stokes} equations and backward
  uniqueness.
\newblock \textit{Russian Math.\ Surveys}, 58(2):211--250, 2003.

\bibitem{Fefferman2000}
C.\ Fefferman.
\newblock Existence and smoothness of the {Navier-Stokes} equation.
\newblock In \textit{Millennium Prize Problems}, pages 57--67.
  Clay Math.\ Inst., Cambridge, MA, 2000.

\bibitem{Ladyzhenskaya1967}
O.\ A.\ Ladyzhenskaya.
\newblock On uniqueness and smoothness of generalized solutions to the
  {Navier-Stokes} equations.
\newblock \textit{Zapiski Nauchn.\ Sem.\ LOMI}, 5:169--185, 1967.

\bibitem{LadyzhenskayaBook1968}
O.\ A.\ Ladyzhenskaya.
\newblock \textit{The Mathematical Theory of Viscous Incompressible Flow}.
\newblock Gordon and Breach, New York, 2nd edition, 1969.

\bibitem{Leray1934}
J.\ Leray.
\newblock Sur le mouvement d'un liquide visqueux emplissant l'espace.
\newblock \textit{Acta Math.}, 63:193--248, 1934.

\bibitem{Prodi1959}
G.\ Prodi.
\newblock Un teorema di unicità per le equazioni di {Navier-Stokes}.
\newblock \textit{Ann.\ Mat.\ Pura Appl.}, 48:173--182, 1959.

\bibitem{Serrin1962}
J.\ Serrin.
\newblock On the interior regularity of weak solutions of the {Navier-Stokes}
  equations.
\newblock \textit{Arch.\ Rational Mech.\ Anal.}, 9:187--195, 1962.

\bibitem{Tao2016}
T.\ Tao.
\newblock Finite time blowup for an averaged three-dimensional {Navier-Stokes}
  equation.
\newblock \textit{J.\ Amer.\ Math.\ Soc.}, 29(3):601--674, 2016.

\bibitem{TFirstPaper3}
P.\ Menon and Claude (Anthropic).
\newblock {The CZ Source Structure Delivers the Delta-Gain: Unconditional
  $H^{1+\delta}$ Regularity for the Navier-Stokes Auxiliary Scalar}.
\newblock \textit{T-First Programme, Paper~3}, 2026.
  In preparation.

\bibitem{TFirstPRD2026}
P.\ Menon and Claude (Anthropic).
\newblock {T-First Programme Primary Reference Document, Version~1.0~FINAL}.
\newblock Technical report, 2026.

\bibitem{Wang2025}
X.\ Wang.
\newblock Self-similar profiles and blow-up scenarios for the {Navier-Stokes}
  equations: counter-flow and {Boussinesq} instabilities.
\newblock Preprint, 2025.

\end{thebibliography}

%======================================================================
%  Appendix
%======================================================================

\appendix

\section{Proof of the Parabolic Schauder Estimate (\cref{lem:parabolic-schauder})}\label{app:parabolic}

We give a self-contained proof of the parabolic estimate used in Step~1.

\begin{proof}[Proof of \cref{lem:parabolic-schauder}]
  Let $\theta$ solve $\pa_t\theta + u\cdot\grad\theta - \nu\lap\theta = S$
  with $\theta(\cdot,0)=0$ and $S\in L^{1+\eps}(\R^3\times[0,T])$.

  \textit{Step A: Homogeneous estimate.}
  By the $L^p$ maximal regularity for the heat equation
  (Da Prato-Grisvard theorem, see \cite{Amann1997}), the solution $\varphi$ of
  $\pa_t\varphi - \nu\lap\varphi = S$, $\varphi(\cdot,0)=0$,
  satisfies
  \[
    \norm{\pa_t\varphi}_{L^{1+\eps}_{t,x}}
    + \norm{D^2\varphi}_{L^{1+\eps}_{t,x}}
    \leq C(\nu,\eps)\,\norm{S}_{L^{1+\eps}_{t,x}}.
  \]
  In particular $\varphi\in W^{1,1+\eps}_t W^{2,1+\eps}_x$.

  \textit{Step B: Transport perturbation.}
  The full equation is $\pa_t\theta - \nu\lap\theta = S - u\cdot\grad\theta$.
  Treat the transport term as a perturbation.
  Since $u\in L^2_t H^1_x\hookrightarrow L^2_t L^6_x$ and
  $\grad\theta\in L^2_t L^2_x$ (from the $H^1$ energy estimate in
  \cref{lem:theta-H1}), the product
  $u\cdot\grad\theta\in L^1_{t,x}$.
  By the fixed-point argument (contraction on $L^{1+\eps}$ in time on small intervals,
  then stitched globally), $\theta$ inherits the $W^{2,1+\eps}$ estimate.

  \textit{Step C: Sobolev embedding.}
  The Sobolev embedding $W^{2,1+\eps}(\R^3)\hookrightarrow H^{1+\delta}(\R^3)$
  holds for $\delta = \delta(\eps) = \eps\cdot 3/(2(2+\eps)) > 0$
  (by the Sobolev embedding theorem applied to the Bessel-potential scale).
  Since $\eps>0$, we have $\delta>0$.

  The estimate
  $\norm{\theta}_{L^2_t H^{1+\delta}_x}\leq C\,\norm{\theta}_{L^2_t W^{2,1+\eps}_x}
  \leq C(\eps,\nu,T)\,\norm{S}_{L^{1+\eps}_{t,x}}$
  follows by combining Steps~A, B, C.
\end{proof}

\section{Numerical Parameters}\label{app:numerics}

\begin{table}[h]
\centering
\caption{Numerical parameters for the experiments reported in \cref{sec:numerics}.}
\label{tab:numerics}
\begin{tabular}{lll}
\toprule
Parameter & 2D value & 3D value \\
\midrule
Grid size $N$         & $64$       & $64^3$     \\
Viscosity $\nu$       & $0.01$     & $0.01$     \\
Final time $T$        & $1.0$      & $1.0$      \\
Time step (adaptive)  & CFL $\leq 0.5$ & CFL $\leq 0.5$ \\
Dealiasing            & 2/3-rule   & 2/3-rule   \\
Time scheme           & RK4        & RK4 + CN (for $\theta$) \\
Hardware              & CPU (NumPy) & JAX Metal (M5)\\
$\theta$ extraction   & H$^s$ norms, $s\in[1,2]$ & same \\
Seed                  & 42         & 42         \\
\bottomrule
\end{tabular}
\end{table}

The $\delta_{\max}$ values in \cref{tab:2d-results,tab:3d-results} are the
maximum observed over all time steps in $[0,T]$.
The extraction method \eqref{eq:delta-extraction} is stable to within
$\pm 0.05$ across repeated runs with different random seeds.

\end{document}
